\documentclass[a4paper,11pt]{article}
\usepackage[utf8]{inputenc}
\usepackage[T1]{fontenc}
\usepackage[french]{babel}
\usepackage{geometry}
\geometry{margin=2.5cm}
\usepackage{hyperref}
\usepackage{booktabs}
\usepackage{longtable}

\title{Journal de bord --- Experimentation CSP benzenoides}
\author{Toky}
\date{Avril 2026 --}

\begin{document}
\maketitle

\tableofcontents
\newpage

%% ===================================================================
\section*{16 avril 2026 --- Mise en place de l'infrastructure}

\textbf{Objectif :} Creer l'architecture d'experimentation (dossiers, scripts, documents).

\textbf{Travail effectue :}
\begin{itemize}
    \item Creation du dossier \texttt{experiments/} avec sous-dossiers organises
    \item Implementation de \texttt{hex\_grid.py} (systeme de coordonnees hexagonales)
    \item Implementation de \texttt{generate\_graph.py} (export .graph)
    \item Implementation de \texttt{shapes.py} (templates de formes)
    \item Implementation de \texttt{benzdb\_import.py} (import depuis BenzDB)
    \item Implementation de \texttt{batch\_run.py} (execution batch)
    \item Creation des squelettes LaTeX (ce journal + rapport)
\end{itemize}

\textbf{Decisions :}
\begin{itemize}
    \item Priorite 1 : valider le modele sur les benzenoides BenzDB (h$\leq$9)
    \item Reference principale : these Varet (2022) pour les familles de structures
\end{itemize}

\textbf{Prochaine etape :} Importer les benzenoides BenzDB et lancer la Serie 0a (h=3..5).

%% ===================================================================
%% Ajouter les prochaines entrees ici :
%%
%% \section*{[date] --- [titre]}
%% \textbf{Objectif :} ...
%% \textbf{Travail effectue :} ...
%% \textbf{Observations :} ...
%% \textbf{Decisions :} ...
%% \textbf{Prochaine etape :} ...

\end{document}
